% ===========================================================================
%  Conclusiones y Recomendaciones
% ===========================================================================
\chapter*{Conclusiones}
\addcontentsline{toc}{chapter}{Conclusiones}
\markboth{Conclusiones}{Conclusiones}

El trabajo alcanz\'o los objetivos propuestos.  Se dise\~n\'o e implement\'o
un lenguaje de dominio espec\'{\i}fico en Common Lisp para la descripci\'on
declarativa de estructuras tabulares de planificaci\'on, y se valid\'o sobre
tres casos reales de la Facultad de Matem\'atica y Computaci\'on de la
Universidad de La Habana mediante un backend que produce archivos \xlsx{}.

Las conclusiones principales son las siguientes.

\begin{enumerate}

  \item \textbf{La separaci\'on entre descripci\'on sem\'antica y backend de
    generaci\'on es alcanzable con un DSL interno en Lisp.}  El sistema
    permite declarar tablas relacionadas, columnas calculadas, cuantificadores
    existenciales y reglas de coloreado condicional sin hacer referencia a
    ning\'un detalle concreto del sistema de salida: ning\'un \'{\i}ndice de
    columna, ninguna cadena de f\'ormula, ninguna llamada a API alguna.  El
    backend es el \'unico punto del sistema donde viven esos detalles, y es
    intercambiable sin modificar el programa DSL.

  \item \textbf{Las macros de Common Lisp son el mecanismo adecuado para
    implementar el DSL.}  La homoiconicidad del lenguaje hace que los macros
    sean transformadores de c\'odigo completamente naturales: cada forma del
    DSL es una lista anidada que el macro convierte en constructores CLOS,
    sin ning\'un parser externo.

  \item \textbf{El despacho m\'ultiple de CLOS ofrece extensibilidad sin
    modificar el c\'odigo existente.}  A\~nadir un nuevo tipo de expresi\'on
    al lenguaje requiere \'unicamente tres acciones: definir la clase del nodo
    con \dsl{defclass*}, definir el macro de superficie y definir el m\'etodo
    del backend.  El patr\'on se replic\'o consistentemente en los m\'as de
    veinte tipos de nodo implementados.

  \item \textbf{El cuantificador \dsl{exists} captura correctamente las
    condiciones de conflicto de planificaci\'on.}  Las dos variantes
    (\dsl{:all-rows} y \dsl{:rows-of}) cubren los patrones de conflicto
    m\'as frecuentes en los documentos analizados.  El backend genera
    f\'ormulas SUMPRODUCT que implementan esas sem\'anticas de forma robusta
    y producen \texttt{FormulaRule} persistentes en el \xlsx{}.

  \item \textbf{El posicionamiento relativo con \dsl{nav} elimina las
    coordenadas absolutas de los programas del DSL.}  Totales, etiquetas y
    celdas de resumen se colocan siempre en la posici\'on correcta
    independientemente del n\'umero de filas de datos, lo que hace al
    programa resiliente ante cambios en el tama\~no de los conjuntos.

  \item \textbf{El backend implementado genera archivos \xlsx{} con formato
    condicional activo.}  La decisi\'on de usar \texttt{FormulaRule} en lugar
    de coloreado est\'atico demuestra que la arquitectura es capaz de producir
    salidas con comportamiento din\'amico: la hoja de c\'alculo sigue siendo
    funcional cuando el usuario la edita directamente, sin necesidad de volver
    a ejecutar el DSL.  Un backend alternativo podr\'{\i}a reproducir esta
    propiedad en su propio formato de salida.

\end{enumerate}

% ---------------------------------------------------------------------------
\chapter*{Recomendaciones}
\addcontentsline{toc}{chapter}{Recomendaciones}
\markboth{Recomendaciones}{Recomendaciones}

A partir del trabajo realizado se identifican las siguientes l\'{\i}neas de
extensi\'on para investigaciones futuras.

\begin{enumerate}

  \item \textbf{Soporte para m\'as funciones de celda.}  El DSL actual cubre
    \texttt{COUNTIF}, \texttt{COUNTA}, \texttt{SUM}, \texttt{IF},
    \texttt{AVERAGEIF} y operaciones de tiempo.  Funciones como
    \texttt{INDEX/MATCH}, \texttt{XLOOKUP}, \texttt{PIVOTBY} y las nuevas
    funciones de matrices din\'amicas de Excel~365 podr\'{\i}an a\~nadirse
    siguiendo exactamente el mismo patr\'on de extensi\'on descrito en la
    secci\'on~\ref{sec:extensibilidad}.

  \item \textbf{Backends alternativos.}  La independencia del backend respecto
    al AST hace posible implementar generadores de c\'odigo para otros
    formatos: ODS (LibreOffice Calc), CSV anotado, HTML con tablas,
    o incluso bases de datos SQL.  La capa del DSL no requiere ning\'un cambio.

  \item \textbf{Interfaz de usuario gr\'afica.}  El AST es una estructura de
    datos de primera clase en Lisp.  Una interfaz gr\'afica podr\'{\i}a visualizar
    y editar el AST directamente, permitiendo a usuarios no programadores
    configurar las reglas de coloreado y los c\'alculos sin escribir c\'odigo.

  \item \textbf{Verif\'{\i}caci\'on de tipos del DSL.}  El sistema actual no
    comprueba en tiempo de expansi\'on de macros que los s\'{\i}mbolos de columna
    referenciados en \dsl{exists} y \dsl{trange} existan en las tablas
    correspondientes.  Un verificador de tipos a nivel del AST detectar\'{\i}a
    estos errores antes de la generaci\'on del c\'odigo Python.

  \item \textbf{Generaci\'on incremental.}  El sistema genera el programa
    Python completo en cada ejecuci\'on.  Un modo incremental que detecte
    qu\'e tablas cambiaron y s\'olo regenere las hojas afectadas reducir\'{\i}a
    el tiempo de ciclo en workbooks con muchas hojas.

  \item \textbf{Plantillas reutilizables.}  Las plantillas de tabla definidas
    con \dsl{def-table} podr\'{\i}an exportarse a una biblioteca de plantillas
    est\'andar para los documentos m\'as comunes de la facultad (horarios de
    docencia, defensas, actas de ex\'amen), reduciendo el esfuerzo de
    configuraci\'on inicial para cada nuevo caso.

\end{enumerate}
